\section{}

\providecommand{\eventanchor}[2]{\hypertarget{evt:#1}{#2}}

\paragraph{960年正月}\eventref{NS-0960-CHENQIAO}{NS-0960-CHENQIAO}　陈桥兵变后赵匡胤即帝位，改国号宋。此项以《宋史》卷1〈太祖本纪一〉；《续资治通鉴长编》建隆元年正月条；《宋会要辑稿·礼》建国礼仪条 为基本来源线索；编年与官修材料能够确认年序和制度用语，但对地方执行、人数、动机或社会后果的判断仍须同文书、碑刻、考古及对方叙事互校。

\paragraph{960年}\eventanchor{NS-0960-KAIFENG-CAPITAL}{NS-0960-KAIFENG-CAPITAL}　宋廷承用开封为东京和中央行政中心。此项以《宋史》卷1〈太祖本纪一〉；《宋会要辑稿·方域》京师条；开封北宋东京城遗址城垣、街道与水系考古资料 为基本来源线索；编年与官修材料能够确认年序和制度用语，但对地方执行、人数、动机或社会后果的判断仍须同文书、碑刻、考古及对方叙事互校。

\paragraph{960年}\eventanchor{NS-0960-REGNAL-YEAR}{NS-0960-REGNAL-YEAR}　宋以宋太祖在位第1年作为诏令、赋役与外交文书的纪年框架。此项以《宋史》卷1—23〈本纪〉；《续资治通鉴长编》本年条；《宋会要辑稿》相关门类 为基本来源线索；编年与官修材料能够确认年序和制度用语，但对地方执行、人数、动机或社会后果的判断仍须同文书、碑刻、考古及对方叙事互校。

\paragraph{961年前后}\eventanchor{NS-0961-MILITARY-POWER}{NS-0961-MILITARY-POWER}　太祖以宴饮、任官和调动方式削弱禁军宿将的独立兵权。此项以《宋史》卷1〈太祖本纪一〉；《宋史》卷250〈石守信传〉；《续资治通鉴长编》建隆二年条 为基本来源线索；编年与官修材料能够确认年序和制度用语，但对地方执行、人数、动机或社会后果的判断仍须同文书、碑刻、考古及对方叙事互校。

\paragraph{961年}\eventanchor{NS-0961-REGNAL-YEAR}{NS-0961-REGNAL-YEAR}　宋以宋太祖在位第2年作为诏令、赋役与外交文书的纪年框架。此项以《宋史》卷1—23〈本纪〉；《续资治通鉴长编》本年条；《宋会要辑稿》相关门类 为基本来源线索；编年与官修材料能够确认年序和制度用语，但对地方执行、人数、动机或社会后果的判断仍须同文书、碑刻、考古及对方叙事互校。

\paragraph{962年}\eventanchor{NS-0962-REGNAL-YEAR}{NS-0962-REGNAL-YEAR}　宋以宋太祖在位第3年作为诏令、赋役与外交文书的纪年框架。此项以《宋史》卷1—23〈本纪〉；《续资治通鉴长编》本年条；《宋会要辑稿》相关门类 为基本来源线索；编年与官修材料能够确认年序和制度用语，但对地方执行、人数、动机或社会后果的判断仍须同文书、碑刻、考古及对方叙事互校。

\paragraph{963年}\eventanchor{NS-0963-JINGNAN}{NS-0963-JINGNAN}　宋军入荆南，结束高氏政权。此项以《宋史》卷2〈太祖本纪二〉；《续资治通鉴长编》乾德元年条；《宋会要辑稿·方域》荆南条 为基本来源线索；编年与官修材料能够确认年序和制度用语，但对地方执行、人数、动机或社会后果的判断仍须同文书、碑刻、考古及对方叙事互校。

\paragraph{963年}\eventanchor{NS-0963-REGNAL-YEAR}{NS-0963-REGNAL-YEAR}　宋以宋太祖在位第4年作为诏令、赋役与外交文书的纪年框架。此项以《宋史》卷1—23〈本纪〉；《续资治通鉴长编》本年条；《宋会要辑稿》相关门类 为基本来源线索；编年与官修材料能够确认年序和制度用语，但对地方执行、人数、动机或社会后果的判断仍须同文书、碑刻、考古及对方叙事互校。

\paragraph{964年冬}\eventanchor{NS-0964-LATER-SHU-CAMPAIGN}{NS-0964-LATER-SHU-CAMPAIGN}　宋军发动伐后蜀战役。此项以《宋史》卷2〈太祖本纪二〉；《宋史》卷479〈后蜀世家〉；《续资治通鉴长编》乾德二年冬条 为基本来源线索；编年与官修材料能够确认年序和制度用语，但对地方执行、人数、动机或社会后果的判断仍须同文书、碑刻、考古及对方叙事互校。

\paragraph{964年}\eventanchor{NS-0964-REGNAL-YEAR}{NS-0964-REGNAL-YEAR}　宋以宋太祖在位第5年作为诏令、赋役与外交文书的纪年框架。此项以《宋史》卷1—23〈本纪〉；《续资治通鉴长编》本年条；《宋会要辑稿》相关门类 为基本来源线索；编年与官修材料能够确认年序和制度用语，但对地方执行、人数、动机或社会后果的判断仍须同文书、碑刻、考古及对方叙事互校。

\paragraph{965年}\eventanchor{NS-0965-LATER-SHU-ANNEXATION}{NS-0965-LATER-SHU-ANNEXATION}　孟昶降宋，后蜀政权终结。此项以《宋史》卷2〈太祖本纪二〉；《宋史》卷479〈后蜀世家〉；《续资治通鉴长编》乾德三年条 为基本来源线索；编年与官修材料能够确认年序和制度用语，但对地方执行、人数、动机或社会后果的判断仍须同文书、碑刻、考古及对方叙事互校。

\paragraph{965年}\eventanchor{NS-0965-REGNAL-YEAR}{NS-0965-REGNAL-YEAR}　宋以宋太祖在位第6年作为诏令、赋役与外交文书的纪年框架。此项以《宋史》卷1—23〈本纪〉；《续资治通鉴长编》本年条；《宋会要辑稿》相关门类 为基本来源线索；编年与官修材料能够确认年序和制度用语，但对地方执行、人数、动机或社会后果的判断仍须同文书、碑刻、考古及对方叙事互校。

\paragraph{966年}\eventanchor{NS-0966-REGNAL-YEAR}{NS-0966-REGNAL-YEAR}　宋以宋太祖在位第7年作为诏令、赋役与外交文书的纪年框架。此项以《宋史》卷1—23〈本纪〉；《续资治通鉴长编》本年条；《宋会要辑稿》相关门类 为基本来源线索；编年与官修材料能够确认年序和制度用语，但对地方执行、人数、动机或社会后果的判断仍须同文书、碑刻、考古及对方叙事互校。

\paragraph{967年}\eventanchor{NS-0967-REGNAL-YEAR}{NS-0967-REGNAL-YEAR}　宋以宋太祖在位第8年作为诏令、赋役与外交文书的纪年框架。此项以《宋史》卷1—23〈本纪〉；《续资治通鉴长编》本年条；《宋会要辑稿》相关门类 为基本来源线索；编年与官修材料能够确认年序和制度用语，但对地方执行、人数、动机或社会后果的判断仍须同文书、碑刻、考古及对方叙事互校。

\paragraph{968年}\eventanchor{NS-0968-REGNAL-YEAR}{NS-0968-REGNAL-YEAR}　宋以宋太祖在位第9年作为诏令、赋役与外交文书的纪年框架。此项以《宋史》卷1—23〈本纪〉；《续资治通鉴长编》本年条；《宋会要辑稿》相关门类 为基本来源线索；编年与官修材料能够确认年序和制度用语，但对地方执行、人数、动机或社会后果的判断仍须同文书、碑刻、考古及对方叙事互校。

\paragraph{969年}\eventanchor{NS-0969-REGNAL-YEAR}{NS-0969-REGNAL-YEAR}　宋以宋太祖在位第10年作为诏令、赋役与外交文书的纪年框架。此项以《宋史》卷1—23〈本纪〉；《续资治通鉴长编》本年条；《宋会要辑稿》相关门类 为基本来源线索；编年与官修材料能够确认年序和制度用语，但对地方执行、人数、动机或社会后果的判断仍须同文书、碑刻、考古及对方叙事互校。

\paragraph{970年}\eventanchor{NS-0970-REGNAL-YEAR}{NS-0970-REGNAL-YEAR}　宋以宋太祖在位第11年作为诏令、赋役与外交文书的纪年框架。此项以《宋史》卷1—23〈本纪〉；《续资治通鉴长编》本年条；《宋会要辑稿》相关门类 为基本来源线索；编年与官修材料能够确认年序和制度用语，但对地方执行、人数、动机或社会后果的判断仍须同文书、碑刻、考古及对方叙事互校。

\paragraph{971年}\eventanchor{NS-0971-REGNAL-YEAR}{NS-0971-REGNAL-YEAR}　宋以宋太祖在位第12年作为诏令、赋役与外交文书的纪年框架。此项以《宋史》卷1—23〈本纪〉；《续资治通鉴长编》本年条；《宋会要辑稿》相关门类 为基本来源线索；编年与官修材料能够确认年序和制度用语，但对地方执行、人数、动机或社会后果的判断仍须同文书、碑刻、考古及对方叙事互校。

\paragraph{971年}\eventanchor{NS-0971-SOUTHERN-HAN}{NS-0971-SOUTHERN-HAN}　宋军攻取广州，南汉刘鋹降服。此项以《宋史》卷2〈太祖本纪二〉；《宋史》卷480〈南汉世家〉；《续资治通鉴长编》开宝四年条 为基本来源线索；编年与官修材料能够确认年序和制度用语，但对地方执行、人数、动机或社会后果的判断仍须同文书、碑刻、考古及对方叙事互校。

\paragraph{972年}\eventanchor{NS-0972-REGNAL-YEAR}{NS-0972-REGNAL-YEAR}　宋以宋太祖在位第13年作为诏令、赋役与外交文书的纪年框架。此项以《宋史》卷1—23〈本纪〉；《续资治通鉴长编》本年条；《宋会要辑稿》相关门类 为基本来源线索；编年与官修材料能够确认年序和制度用语，但对地方执行、人数、动机或社会后果的判断仍须同文书、碑刻、考古及对方叙事互校。

\paragraph{973年}\eventanchor{NS-0973-REGNAL-YEAR}{NS-0973-REGNAL-YEAR}　宋以宋太祖在位第14年作为诏令、赋役与外交文书的纪年框架。此项以《宋史》卷1—23〈本纪〉；《续资治通鉴长编》本年条；《宋会要辑稿》相关门类 为基本来源线索；编年与官修材料能够确认年序和制度用语，但对地方执行、人数、动机或社会后果的判断仍须同文书、碑刻、考古及对方叙事互校。

\paragraph{974年}\eventanchor{NS-0974-REGNAL-YEAR}{NS-0974-REGNAL-YEAR}　宋以宋太祖在位第15年作为诏令、赋役与外交文书的纪年框架。此项以《宋史》卷1—23〈本纪〉；《续资治通鉴长编》本年条；《宋会要辑稿》相关门类 为基本来源线索；编年与官修材料能够确认年序和制度用语，但对地方执行、人数、动机或社会后果的判断仍须同文书、碑刻、考古及对方叙事互校。

\paragraph{974年秋}\eventanchor{NS-0974-SOUTHERN-TANG-CAMPAIGN}{NS-0974-SOUTHERN-TANG-CAMPAIGN}　宋军渡江围攻金陵，展开灭南唐战争。此项以《宋史》卷2〈太祖本纪二〉；《宋史》卷478〈南唐世家〉；《续资治通鉴长编》开宝七年条 为基本来源线索；编年与官修材料能够确认年序和制度用语，但对地方执行、人数、动机或社会后果的判断仍须同文书、碑刻、考古及对方叙事互校。

\paragraph{975年}\eventanchor{NS-0975-REGNAL-YEAR}{NS-0975-REGNAL-YEAR}　宋以宋太祖在位第16年作为诏令、赋役与外交文书的纪年框架。此项以《宋史》卷1—23〈本纪〉；《续资治通鉴长编》本年条；《宋会要辑稿》相关门类 为基本来源线索；编年与官修材料能够确认年序和制度用语，但对地方执行、人数、动机或社会后果的判断仍须同文书、碑刻、考古及对方叙事互校。

\paragraph{975年十一月}\eventanchor{NS-0975-SOUTHERN-TANG-ANNEXATION}{NS-0975-SOUTHERN-TANG-ANNEXATION}　金陵陷落，李煜降宋，南唐灭亡。此项以《宋史》卷2〈太祖本纪二〉；《宋史》卷478〈南唐世家〉；《续资治通鉴长编》开宝八年十一月条 为基本来源线索；编年与官修材料能够确认年序和制度用语，但对地方执行、人数、动机或社会后果的判断仍须同文书、碑刻、考古及对方叙事互校。

\paragraph{976年}\eventanchor{NS-0976-REGNAL-YEAR}{NS-0976-REGNAL-YEAR}　宋以宋太祖在位第17年作为诏令、赋役与外交文书的纪年框架。此项以《宋史》卷1—23〈本纪〉；《续资治通鉴长编》本年条；《宋会要辑稿》相关门类 为基本来源线索；编年与官修材料能够确认年序和制度用语，但对地方执行、人数、动机或社会后果的判断仍须同文书、碑刻、考古及对方叙事互校。

\paragraph{976年十月}\eventanchor{NS-0976-TAIZU-SUCCESSION}{NS-0976-TAIZU-SUCCESSION}　太祖去世，赵光义即位为太宗。此项以《宋史》卷2〈太祖本纪二〉；《宋史》卷3〈太宗本纪一〉；《续资治通鉴长编》开宝九年十月条 为基本来源线索；编年与官修材料能够确认年序和制度用语，但对地方执行、人数、动机或社会后果的判断仍须同文书、碑刻、考古及对方叙事互校。

\paragraph{977年}\eventanchor{NS-0977-REGNAL-YEAR}{NS-0977-REGNAL-YEAR}　宋以宋太宗在位第1年作为诏令、赋役与外交文书的纪年框架。此项以《宋史》卷1—23〈本纪〉；《续资治通鉴长编》本年条；《宋会要辑稿》相关门类 为基本来源线索；编年与官修材料能够确认年序和制度用语，但对地方执行、人数、动机或社会后果的判断仍须同文书、碑刻、考古及对方叙事互校。

\paragraph{978年}\eventanchor{NS-0978-REGNAL-YEAR}{NS-0978-REGNAL-YEAR}　宋以宋太宗在位第2年作为诏令、赋役与外交文书的纪年框架。此项以《宋史》卷1—23〈本纪〉；《续资治通鉴长编》本年条；《宋会要辑稿》相关门类 为基本来源线索；编年与官修材料能够确认年序和制度用语，但对地方执行、人数、动机或社会后果的判断仍须同文书、碑刻、考古及对方叙事互校。

\paragraph{979年六月}\eventanchor{NS-0979-GAOLIANG-RIVER}{NS-0979-GAOLIANG-RIVER}　宋军北攻幽州，在高梁河失利后南撤。此项以《宋史》卷3〈太宗本纪一〉；《辽史》〈景宗本纪〉；《续资治通鉴长编》太平兴国四年六月条 为基本来源线索；编年与官修材料能够确认年序和制度用语，但对地方执行、人数、动机或社会后果的判断仍须同文书、碑刻、考古及对方叙事互校。

\paragraph{979年五月}\eventanchor{NS-0979-NORTHERN-HAN}{NS-0979-NORTHERN-HAN}　宋灭北汉，完成对中原和河东割据政权的统一。此项以《宋史》卷3〈太宗本纪一〉；《宋史》卷480〈北汉世家〉；《续资治通鉴长编》太平兴国四年五月条 为基本来源线索；编年与官修材料能够确认年序和制度用语，但对地方执行、人数、动机或社会后果的判断仍须同文书、碑刻、考古及对方叙事互校。

\paragraph{979年}\eventanchor{NS-0979-REGNAL-YEAR}{NS-0979-REGNAL-YEAR}　宋以宋太宗在位第3年作为诏令、赋役与外交文书的纪年框架。此项以《宋史》卷1—23〈本纪〉；《续资治通鉴长编》本年条；《宋会要辑稿》相关门类 为基本来源线索；编年与官修材料能够确认年序和制度用语，但对地方执行、人数、动机或社会后果的判断仍须同文书、碑刻、考古及对方叙事互校。

\paragraph{980年}\eventanchor{NS-0980-REGNAL-YEAR}{NS-0980-REGNAL-YEAR}　宋以宋太宗在位第4年作为诏令、赋役与外交文书的纪年框架。此项以《宋史》卷1—23〈本纪〉；《续资治通鉴长编》本年条；《宋会要辑稿》相关门类 为基本来源线索；编年与官修材料能够确认年序和制度用语，但对地方执行、人数、动机或社会后果的判断仍须同文书、碑刻、考古及对方叙事互校。

\paragraph{981年}\eventanchor{NS-0981-REGNAL-YEAR}{NS-0981-REGNAL-YEAR}　宋以宋太宗在位第5年作为诏令、赋役与外交文书的纪年框架。此项以《宋史》卷1—23〈本纪〉；《续资治通鉴长编》本年条；《宋会要辑稿》相关门类 为基本来源线索；编年与官修材料能够确认年序和制度用语，但对地方执行、人数、动机或社会后果的判断仍须同文书、碑刻、考古及对方叙事互校。

\paragraph{982年}\eventanchor{NS-0982-REGNAL-YEAR}{NS-0982-REGNAL-YEAR}　宋以宋太宗在位第6年作为诏令、赋役与外交文书的纪年框架。此项以《宋史》卷1—23〈本纪〉；《续资治通鉴长编》本年条；《宋会要辑稿》相关门类 为基本来源线索；编年与官修材料能够确认年序和制度用语，但对地方执行、人数、动机或社会后果的判断仍须同文书、碑刻、考古及对方叙事互校。

